\section{Discussion}
\label{sec:discussion}

\IC{} carried one demonstration from a test system through posteriors to predictions. On M0, both inference routes pass simulation-based calibration, with small
declared discrepancies, and the posterior-to-prediction step covers lactate export at the expected rate, though the data measure export noisily through extracellular lactate, and how far that drives the posterior is unmeasured. The knockdown's effect is what the data did not see. On the reduced model the predictions are unscored, and their budget depends strongly on
a prior width that we chose. Any user can release fixed constants to their priors, and releasing \pn{traj.nreleased} of them raises the predicted mean GTP at the end of the cycle about fivefold, unscored and on one posterior (\cref{sec:trajectories}). The M0 check does not test the completion or the transfer into the reduced model.

\paragraph{What scales and what does not.} Exact inference is calibrated on M0, and it stops there. The cost of an exact posterior grows with
the number of genes and the length of the record, and the full reduced model is
already above our budget. The mixed route removes the stochastic paths from
the sampler. Two quantities limit it. The first is the Jensen gap, which we declare
and correct (\cref{sec:jensen}). The second is $\mathcal{R}$, the variance of the bulk log-likelihood over the stochastic paths. As an upper bound, at nominal parameters, $\mathcal{R}$ is about \pn{info.Rmzero} on M0 at \pn{mzero.cells} cells and \pn{info.Rcorea} on the full reduced model at 200 cells, both at 10\% assay noise. If $\mathcal{R}$ is near those bounds, the metabolites constrain the paths together, and a single draw of paths per parameter value is not enough. On M0, the Gaussian bulk factor passes calibration
in this regime (\cref{sec:posteriors}). We have not tested it on the full reduced model.

\paragraph{The observation model sets the interface.} The three graphs of
\cref{sec:info} suggest a working rule, from one model and two directions. Single-cell data inform the parameters of their
own module. Information between modules travels on bulk observables, and it follows
the gain of the coupling. A modeller can compute the gains before any inference. In this work they anticipated the transcript null. For the other direction the evidence is thinner: on one dataset the mixed route shrinks the ptsG promoter by less than the linearised bound projected, and the exact sampler did not measure it.

\paragraph{Reductions as approximations to inference.} A dynamical reduction changes
the simulator and keeps full Bayesian inference. A statistical cut keeps the
simulator and changes the target. We prefer the first, because its cost can be
measured against a more detailed composition. The formalism swap is such a reduction.
We do not claim that the more detailed composition is the truth about the cell. We claim that each choice of description is declared, and that the swap's cost is measured on M0.

The prediction of \cref{sec:budget} uses the other tool, a cut, and it does so on
purpose. There the choice is not between two ways of inferring one model. No data
reach the held parameters at all, and an exact posterior on the reduced model is out of
reach. So the cut says plainly what the prediction is: a modular transfer of an M0
posterior, with the other parameters at their priors. Its cost is not measured, and it
is declared (\cref{tab:modelform}).

\subsection{Limitations}

\begin{itemize}
  \item \emph{Synthetic data only.} The data come from the model itself. Real
  assays need observation models for unknown cell numbers in bulk metabolomics and
  for capture efficiency in single-cell RNA counts.
  \item \emph{One organism, one reduction.} We demonstrate the framework on reduced
  \syn{} only, and it has no generic organism registry. Generality across organisms needs
  a second implementation.
  \item \emph{Asserted priors.} Many priors have no width in their source, so we set
  the width. This choice matters: the held share of the budget runs from
  \pn{budget.heldpctlo}\% to \pn{budget.heldpcthi}\% across the widths we tried
  (\cref{sec:budget}).
  \item \emph{Exact calibration at small scale only.} Coverage of the exact sampler
  at the full reduced model is out of reach at about \pn{cost.fullscale} CPU-hours per
  posterior. The reduced model's prediction is therefore unscored, and it rests on one
  posterior and on a transfer we have not tested.
  \item \emph{Weak reverse coupling.} The mixed route relies on the small
  elasticities of the rebuild (\pn{gain.txlo}--\pn{gain.tl}). On the reduced model about \pn{kd.hi}\% of baseline cells already fall below $K_d$ at least once (\cref{sec:modelform}), so this basis is in question there without any perturbation.
\end{itemize}

\subsection{Future work}

Four directions follow. The first is the method's own limits: conditioning the bulk factor on the observed transcripts, sampling paths only for pools near the one-particle floor, and handling strong-feedback regimes. The second is scale and data: the full well-stirred
\syn{} model through the fast port of the companion paper \cite{kimpson2026whole}, then
real measurements. The third is comparing model structure and formalisms, which the swap
makes possible, by predictive scores rather than Bayes factors. The fourth is goal-directed
inference and experimental design: which parts of a coupled simulator must be kept, and
which latent variables inferred, to keep the information a prediction needs?
